\documentclass[instructions.tex]{subfiles}
\begin{document}

\chapter{Autres motifs de conversion, tirés de la passion de notre Seigneur Jésus-Christ.}

CONVERTISSEZ-VOUS pour l'amour de Jésus-Christ. Il le mérite ; il le demande.

\primo Il le mérite, 1. par lui-même. Il est votre roi et votre maître, votre juge et votre rémunérateur, votre sauveur et votre Dieu. Quels sentimens de fidélité et d'obéissance, de respect et de reconnoissance, de confiance et d'amour ne lui devez-vous pas pour ces qualités ?

2. Il le mérite par les tourmens qu'il a soufferts. Pouvez-vous y penser sans avoir le cœur attendri ? Un animal privé de raison a du retour pour son maître, il le défend jusqu'à se laisser charger de coups, pour une vile nourriture qu'il en reçoit ; et nous, pour qui Jésus-Christ a donné sa vie, nous ne lui rendons que des ingratitudes.

Nous avons de la reconnoissance envers nos amis pour les moindres services. Jésus-Christ sera-t-il le seul ami pour lequel nous n'aurons que de la dureté ? Qui de nos amis, après tout, a été crucifié pour nous ? Si un ennemi avait exposé sa vie pour nous secourir, nous lui témoignerions du retour. Traiterons-nous Jésus-Christ plus mal qu'un ennemi ? Quand il n'aurait versé qu'une larme pour nous sauver, notre vie devrait se passer dans des transports de reconnoissance. Que lui devons-nous donc pour tant de sang répandu ?

3. Il mérite notre retour par le sujet qui l'a fait mourir. N'accusons pas tant les juifs de sa mort que nous-mêmes ; les bourreaux ont été les instrumens de nos péchés ; je dois en particulier me reconnoître aussi coupable de sa mort que si j'étois le seul qui l'eusse fait mourir ; je lui suis aussi redevable que s'il étoit mort pour moi seul. Il est vrai que tous les hommes ont contribué à sa mort par leurs péchés ; mais il n'est pas moins vrai que, si je n'avois pas tant péché, il n'auroit pas tant souffert, et que l'affliction qu'il a eue de mes péchés seuls, étoit capable de lui causer la mort.

Un Dieu qui souffre et qui meurt ! voilà, ô mon âme ! l'effet de tes péchés. Un Dieu crucifié, couvert de sang et de plaies ! voilà l'ouvrage de vos mains. O mon Sauveur ! comment puis-je encore aimer le péché, en voyant ce qu'il vous a fait souffrir ? O hommes ! comprenez ce que vous avez coûté à un Dieu, puisque, ayant d'une seule parole formé l'univers, il faut qu'il sue le sang et qu'il souffre la mort pour expier vos péchés.

4. Enfin, Jésus-Christ mérite du retour pour la charité qu'il a eue pour nous. Il n'y a point de charité plus grande que de donner sa vie pour ses amis ; mais il a poussé la charité plus loin, puisqu'il a donné sa vie pour ses ennemis. Oui, nous étions tous ennemis de Dieu et enfans de sa colère. Oh ! quelle miséricorde de mourir ainsi pour des misérables ! O Jésus ! qu'il vous en a coûté de m'aimer, et qu'il m'en coûte peu de vous déplaire !

La charité du Sauveur pour les hommes a été si ardente et si étendue, qu'il n'en est pas un seul pour lequel il n'ait pas offert sur la croix le prix de sa mort à Dieu son Père ; qu'il auroit été disposé, dit saint Jean Chrysostôme, à mourir pour une seule âme, si elle avoit été seule au monde. C'est pourquoi saint Paul, comme si Jésus-Christ étoit mort pour lui seul, ne dit pas seulement qu'il s'est livré pour tous, mais il m'a aimé, et il s'est livré pour moi\footnote{Dilexit me, et tradidit semetipsum pro me. Gal. 2.}. Ah ! dit saint Bernard, si nous ne pouvons faire pour notre Sauveur autant qu'il a fait pour nous, du moins ne l'offensons plus. Il nous a aimés le premier : rendons-lui du moins amour pour amour\footnote{Si amare piget, redamare non piget. S. Bern.}. Jésus-Christ mérite donc notre retour.

\secundo Il le demande, et c'est l'unique chose qu'il vous demande. Il ne pouvoit faire plus qu'il n'a fait pour vous, et il ne peut vous demander moins qu'il ne vous demande : il ne désire que votre amour. Mon fils, donnez-moi votre cœur, cessez de m'offenser\footnote{Præbe fili mi, cor tuum mihi. Prov. 23.}. Le refuserez-vous à ma tendresse, à mes larmes et à mon sang ?

Prenez en main votre crucifix, considérez votre Sauveur ; il expire les bras étendus, son côté ouvert, son cœur percé pour vous y recevoir ; il jette sur vous, du haut de cette croix, ses yeux mourans : Mon fils, je meurs pour toi ; s'il falloit encore souffrir mille fois la mort, je la souffrirois\footnote{Quid debui ultra facere, et non feci ? Is. 5.}.

Tu vois mon corps innocent, déchiré de coups et tout sanglant ; tu me vois expirer dans l'ignominie et dans les douleurs ; jamais il n'y eut de douleur semblable à la mienne ; mais mon plus grand tourment, c'est le péché qui dure, c'est de souffrir pour des ingrats, de mourir pour te sauver, tandis que tu veux te perdre. Je suis accablé par ma douleur ; veux-tu encore l'augmenter en perdant ton âme, pour laquelle j'expire, et qui m'est plus chère que la vie ?

Notre cœur, fût-il aussi dur que les rochers qui se brisèrent à sa mort, pourroit-il résister à des reproches si tendres ? Pousserons-nous l'ingratitude jusqu'à renouveler les souffrances d'un Dieu qui meurt pour nous ? Ah ! puisse la vue de ses larmes et de son sang arrêter le cours de nos péchés, et nous inspirer un amour si ardent, que nous puissions dire, comme saint Paul, que nous ne voulons plus être que pour celui qui est mort pour nous !

Oui, c'est de tout mon cœur, ô Jésus ! que je me convertis à vous ; recevez-moi, pardonnez-moi. Vous êtes venu chercher les pécheurs, les brebis égarées ; je suis, hélas ! plus égaré et plus pécheur que qui que ce soit au monde. C'est parce que je suis misérable, que j'implore votre miséricorde. Souvenez-vous que c'est pour moi que vous avez souffert la mort et répandu votre sang. Faites, ô mon Sauveur et mon Dieu ! que tant de travaux ne soient pas inutiles\footnote{Si in viridi ligno hæc fiunt, in arido quid fiet ? Luc. 23.}.

\section{Résolutions}

\primo Je méditerai, au moins une fois dans la semaine, sur les souffrances de Jésus-Christ dans sa passion.

\secundo Je me considérerai comme cause réelle de ces horribles souffrances.

\tertio Et je regarderai le Fils de Dieu comme substitué à ma place dans le jardin des Oliviers, dans les tribunaux et sur la croix, pour m'arracher à l'enfer, m'apprendre à faire pénitence, et rendre celle que je ferai, efficace. 

\prayer{O mon Sauveur ! prosterné au pied de votre croix, j'y confesse avec confusion mon ingratitude et mes perfidies. Ah ! pardonnez-moi ; c'en est fait, je vous donne mon cœur avec toutes ses affections, crucifiez-le avec vous.}

\end{document}
